\documentclass[10pt,a4paper]{amsart}
\usepackage[margin=19mm]{geometry}
\usepackage{amsmath,amssymb,mathtools}
\usepackage[scr=rsfs,cal=euler]{mathalfa}
\usepackage{xfrac}
\usepackage[unicode,hidelinks,bookmarks=false]{hyperref}
\newcommand{\ie}{i.e.\ }
\newcommand{\cf}{cf.\ }
\newcommand{\resp}{respectively\ }
\newcommand{\Subsection}{\subsection}
\providecommand{\nobreakdash}{\nobreak\hbox{-}}
\setlength{\emergencystretch}{2em}
\multlinegap0pt
\usepackage{amssymb}
\usepackage{xcolor}
\usepackage{algorithm}\usepackage{algpseudocode}
\newtheorem{proposition}{Proposition}
\title{Regional factors determine the feasibility of an elimination strategy}
\author{George Adu-Boahen, Troy Day, Michael J Plank, Amy Hurford}
\date{Source: arXiv:2608.08956v1 (CC BY 4.0)}
\begin{document}
\maketitle

\noindent\emph{Source and licence.} Regional factors determine the feasibility of an elimination strategy. Version arXiv:2608.08956v1 (CC BY 4.0), distributed by arXiv under the Creative Commons Attribution 4.0 International licence, http://creativecommons.org/licenses/by/4.0/, as recorded in the arXiv licence metadata for this exact version and shown on the abstract page https://arxiv.org/abs/2608.08956v1. Full licence notice: \texttt{licenses/arxiv-2608.08956-CC-BY-4.0.txt}.\par\smallskip

\noindent\textit{Editorial scope.} Counted unit: the article's proposition prop:u1-one-switch (source0.tex lines 370--374). The article's complete original proof of it is delivered with it (source0.tex lines 376--450, one \texttt{proof} environment of 75 lines). No mathematics is written, repaired or re-derived: the statement and the proof are the article's own lines, in the article's own notation, and no retained line is substituted or re-flowed.  No further source-local context is needed: every cross-reference of every retained line resolves inside the two delivered spans.  Nothing was omitted from any retained span, so the package carries no omission record.  No retained line contains a citation, so the wrapper carries no bibliography and no bibliography entry.  No retained line contains a figure environment, an image-inclusion command or drawing code, so no image is delivered and the package declares no asset dependency.  Editorial formatting only, in addition: an AMS \texttt{amsart} preamble that restates the article's own declarations and macros used by the retained lines, namely the environments \texttt{proposition} on separate counters, together with the standard packages those lines need.  The retained environments print in this document as Proposition 1 in order of appearance, whereas the article prints the same items with the numbering of its own full document.  The complete native source of the article as distributed in the arXiv e-print for this exact version is bundled unchanged as \texttt{sources/207205/source0.tex}.
\begin{proposition}[The control $u_1$]
\label{prop:u1-one-switch}
The switching function \(\varphi_1\) is strictly decreasing. Consequently,
\(u_1^*\) has at most one switch, necessarily from \(u_{1\max}\) to \(0\).
\end{proposition}

\begin{proof}:
The maximum condition gives
\[
u_1^*(t)
=
\begin{cases}
u_{1\max}, & \varphi_1(t)>0,\\
0, & \varphi_1(t)<0,
\end{cases}
\qquad\text{for a.e. }t.
\]
At points where \(\varphi_1=0\), the value of \(u_1^*\) is not determined by \(\varphi_1\) but the product \(\varphi_1u_1^*\) is always given by \begin{equation}
\label{eq:phi1-u1-product}
\varphi_1u_1^*
=
u_{1\max}\max\{\varphi_1,0\}
\qquad\text{a.e.},
\end{equation}
Differentiating $\varphi_1$ gives
\begin{align}
\dot\varphi_1
&=
-\dot\lambda_{I_1}\notag\\
&=
-\beta\, S\, \Psi-\mu(\lambda_{U_1}^*-\varphi_1)+\varphi_1u_1^*\notag\\
&=
-\beta\, S\, \Psi-\mu \lambda_{U_1}^*+\mu\varphi_1
+u_{1\max}\max\{\varphi_1,0\}.
\label{eq:phi1-first-order}
\end{align}
The right-hand side of \eqref{eq:phi1-first-order} is continuous and therefore \(\varphi_1\in C^1([0,T])\). Moreover, we can differentiate \eqref{eq:phi1-first-order} almost everywhere to obtain
\begin{align}
\ddot \varphi_1
&=
-\frac{d}{dt}[\beta\, S\, \Psi]+(\mu+u_{1\max}\mathbf{1}_{\{\varphi_1>0\}})\dot \varphi_1\notag\\
&=
\left[
\mu
+u_{1\max}\mathbf 1_{\{\varphi_1>0\}}
-\beta S
\right]\dot \varphi_1
\qquad\text{a.e.}
\label{eq:w1-linear}
\end{align}
Define
\[
A(t)
:=
\mu
+u_{1\max}\mathbf 1_{\{\varphi_1(t)>0\}}
-\beta S(t).
\]
Since \(A\in L^\infty(0,T)\), the scalar linear equation
\eqref{eq:w1-linear} can be written
\begin{equation}
\label{eq:w1-representation}
\dot \varphi_1(t)
=
\dot \varphi_1(T)
\exp\left(
-\int_t^T A(s)\,ds
\right).
\end{equation}
By transversality $\lambda_{I_1}(T)=0$ and hence $\varphi_1(T)=\lambda_{U_1}^*\leq 0$. Also, $\beta S(T)\Psi(T)=\beta S(T)$. Therefore, evaluating \eqref{eq:phi1-first-order} at \(T\) gives $\dot \varphi_1(T)=-\beta S(T)<0$. It follows from \eqref{eq:w1-representation} that $\dot\varphi_1(t)<0$ for all $t\in[0,T]$ (a.e.). Thus \(\varphi_1\) is strictly decreasing and can vanish at most once. Since the maximum condition selects \(u_{1\max}\) when
\(\varphi_1>0\) and \(0\) when \(\varphi_1<0\), there exists
\(\tau_1\in[0,T]\) such that, up to equality almost everywhere,
\[
u_1^*(t)
=
\begin{cases}
u_{1\max}, & 0\leq t<\tau_1,\\
0, & \tau_1<t\leq T.
\end{cases}
\]
\end{proof}

\end{document}
